\documentclass[11pt]{article}
\usepackage{fontspec}
\usepackage[french,english]{babel}
\usepackage{amsmath}
\usepackage{amssymb}
\usepackage{amsthm}
\usepackage{mathrsfs}
\usepackage{tikz-cd}
\usepackage{enumitem}

\newcommand{\printedpage}[1]{\marginpar{\footnotesize p.~#1}}
\newcommand{\msp}[1]{\marginpar{\footnotesize [#1]}}
\newcommand{\uTy}{\underline{\mathrm{Ty}}}

\begin{document}

\printedpage{34}
% head: A. Grothendieck : Dérivateurs, Ch. VII, Édition des Universités de Montpellier II et Paris VII
We are going to describe these monoids $\mathrm{Nerf}_i(\uTy)$, and observe that they are all monoids that are \textit{free,} just like $\mathrm{Nerf}_0(\uTy) = \mathrm{Ob}\,\uTy = \mathrm{Ty}$. The semi-simplicial maps between these monoids are therefore known when we know them on the \textit{generators} of these free monoids (which generators are also the indecomposable elements, hence defined intrinsically in terms of the multiplicative structure). % [sic] connus -> connues, agreement with « applications »

For $\mathrm{Nerf}_0(\uTy) = \mathrm{Ty}$, this is something we have seen, the generators are
\[
\Delta_+,\Delta_-,
\]
corresponding respectively to the oriented intervals
\[
\begin{array}{rccc}
\boldsymbol{\Delta}_+ : & 0 & \longrightarrow & 1\\[4pt]
\boldsymbol{\Delta}_- : & 0 & \longleftarrow & 1.
\end{array}
\]

\msp{page 50$'$}
For $\mathrm{Nerf}_1(\uTy) = \mathrm{Fl}(\uTy)$, i.e. for the arrows, result $3^o$) of the preceding paragraph (determination of $\mathrm{Fl}(\uTy)$) is also interpreted by saying that $\mathrm{Fl}(\uTy)$ is the free monoid generated by an infinite system of generators
\[
\mathfrak{f}_{\varepsilon,n}\ , \qquad\quad (\varepsilon,n)\in\{+1,-1\}\times\mathbf{N},
\]
indexed by the product set $\{+1,-1\}\times\mathbf{N}$. Here $\mathfrak{f}_{\varepsilon,n}$ denotes the unique arrow in $\uTy$
\[
\mathfrak{f}_{\varepsilon,l} : \Delta_\varepsilon \longrightarrow \Delta^l_\varepsilon
\]
(where of course the exponent $l$ denotes an exponentiation for the composition law in $\uTy$
\[
\Delta^l_\varepsilon = \underbrace{\Delta_\varepsilon\circ\Delta_\varepsilon\circ\ \cdots\ \circ\Delta_\varepsilon}_{l\text{ copies}},
\]
and not the Cartesian power in Cat, where in $\uTy$ (where it exists only if $l\in\{0,1\}$)). % [sic] où dans -> ou dans, conjunction intended
In the description of $\mathrm{Fl}(\uTy)$ in the preceding paragraph, we wrote the ``general term'' of the free monoid by grouping together the groups of successive functors of the form $f_{+,0}$, $f_{-,0}$, i.e. corresponding to the two indecomposable degeneracy arrows (for $*$) % [sic] corréspondant -> correspondant, accent
\[
\left\{
\begin{array}{lccc}
f_{+,0} : & \Delta_+ & \longrightarrow & e\\[4pt]
f_{-,0} : & \Delta_- & \longrightarrow & e
\end{array}
\right.
\]
(\footnote[15]{\textbf{NB} The empty ``word'' of $\mathrm{Fl}(\uTy)$ corresponds to the arrow $\mathrm{id}_1$, where $1$ is the unit element of Ty, realizing the one-point ``interval''.}). The free submonoid of $\mathrm{Fl}(\uTy)$ generated by these two elements is the monoid of arbitrary (pure) degeneracy arrows
% foot: Version 22 décembre 2001
\printedpage{35}
% head: A. Grothendieck : Dérivateurs, Ch. VII, Édition des Universités de Montpellier II et Paris VII
\msp{page 51}
in $\uTy$:
\[
f_\tau : I_\tau \longrightarrow e \qquad\quad (\tau\in\mathrm{Ty}),
\]
and we have
\[
f_\tau * f_{\tau'} = f_{\tau\tau'},
\]
i.e. the map
\[
\tau \longmapsto f_\tau :\quad \underbrace{\mathrm{Ob}\,(\uTy)}_{=\,\mathrm{Ty}\,=\,\mathrm{Nerf}_0(\uTy)} \quad\longrightarrow\quad \underbrace{\mathrm{Fl}(\uTy)}_{\mathrm{Nerf}_1(\uTy)}
\]
is a monoid homomorphism, taking $\Delta_\varepsilon$ ($\varepsilon\in\{+1,-1\}$) to $\mathfrak{f}_{\varepsilon,0}$. This is not the degeneracy map from $\mathrm{Nerf}_0$ to $\mathrm{Nerf}_1$, i.e. the map $\tau\longmapsto\mathrm{id}_\tau$. The latter is given by % [sic] dégénérenscence -> dégénérescence, spelling
\[
\begin{array}{ccc}
\mathrm{Ob}\,(\uTy) & \longrightarrow & \mathrm{Fl}(\uTy)\\[4pt]
\tau & \longmapsto & \mathrm{id}_\tau
\end{array}
\ :\
\left|\
\begin{array}{rcccl}
 & \Delta_\varepsilon & \longmapsto & \mathfrak{f}_{\varepsilon,1} & \qquad(\varepsilon\in\{+1,-1\})\\[4pt]
\text{i.e.} & \Delta_+ & \longmapsto & \mathfrak{f}_{+,1} & \\[4pt]
 & \Delta_- & \longmapsto & \mathfrak{f}_{-,1}\ . &
\end{array}
\right.
\]
As for the source and target homomorphisms
\[
\begin{tikzcd}[column sep=large]
\mathrm{Fl}(\uTy) \arrow[r, shift left, "\sigma\text{ (source)}"] \arrow[r, shift right, "\beta\text{ (target)}"'] & \mathrm{Ob}\,(\uTy),
\end{tikzcd}
\]
they are obvious on the generators
\[
\left\{
\begin{array}{lcl}
\sigma(\mathfrak{f}_{\varepsilon,l}) & = & \Delta_\varepsilon\\[4pt]
\beta(\mathfrak{f}_{\varepsilon,l}) & = & \Delta^l_\varepsilon.
\end{array}
\right.
\]

I am now going to describe $\mathrm{Nerf}_2(\uTy)$, i.e. the systems
\[
\tau \xrightarrow{\ f\ } \tau' \xrightarrow{\ g\ } \tau''.
\]

\msp{page 52}
Giving $f\in\mathrm{Fl}(\uTy)$ amounts to giving an ordered sequence of generators $f_{\varepsilon,l}$, hence a sequence
\[
\begin{aligned}
&\mathfrak{f}_{\varepsilon_1,l_1},\ \mathfrak{f}_{\varepsilon_2,l_2},\ \ \dots\ \ ,\mathfrak{f}_{\varepsilon_n,l_n}\\
&((\varepsilon_i,l_i))_{1\leq i\leq n}\in[\{\pm1\}\times\mathbf{N}]^n.
\end{aligned}
\]
We then have
\[
\begin{aligned}
x &= \Delta_{\varepsilon_1}\Delta_{\varepsilon_2}\cdots\Delta_{\varepsilon_n} &\qquad\qquad& \mathrm{long}\,x = n\\
y &= \Delta^{l_1}_{\varepsilon_1}\Delta^{l_2}_{\varepsilon_2}\cdots\Delta^{l_n}_{\varepsilon_n} && \mathrm{long}\,y = \textstyle\sum_{i\in[1,n]} l_i.
\end{aligned}
\]
% foot: Version 22 décembre 2001

\printedpage{36}
% head: A. Grothendieck : Dérivateurs, Ch. VII, Édition des Universités de Montpellier II et Paris VII
Giving $q : y\longrightarrow z$, i.e. an object of $y\backslash\uTy$, amounts to giving a map from the set of the $m=\sum l_i$ transition arrows of $y$ to $\mathbf{N}$. In fact, we have % [sic] q -> g, the map is called g below (« la donnée de g »)
\[
\mathrm{Tr}(y) = \mathrm{Tr}(\Delta^{l_1}_{\varepsilon_1})\amalg\mathrm{Tr}(\Delta^{l_2}_{\varepsilon_2})\amalg\ \cdots\ \amalg\mathrm{Tr}(\Delta^{l_n}_{\varepsilon_n})
\]
(``sum set'', i.e. disjoint union), hence
\[
\mathrm{Hom}_{\mathrm{Ens}}(\mathrm{Tr}(y)) \simeq \prod_{i\in[1,n]}\mathrm{Hom}_{\mathrm{Ens}}(\underbrace{\mathrm{Tr}(\Delta^{l_i}_{\varepsilon_i})}_{[1,l_i]},\mathbf{N}), % [sic] Hom_Ens(Tr(y)) -> Hom_Ens(Tr(y),N), second argument missing
\]
hence giving $g$ is equivalent to giving a system of $n$ maps
\[
\left\{
\begin{array}{l}
\pi_i : [1,l_i]\longrightarrow\mathbf{N} \qquad\quad (\text{weighting of }\Delta^{l_i}_{\varepsilon_i})\\[6pt]
\text{for } i\in[1,n].
\end{array}
\right.
\]
But for $i\in[1,n]$, let us consider the element of $\mathrm{Nerf}_2$
\[
\Delta_{\varepsilon_1}\xrightarrow{\mathfrak{f}_{\varepsilon_i,l_i}}\Delta^{l_i}_{\varepsilon_i}\xrightarrow{g_{\varepsilon,\pi_i}}\underbrace{\Delta^{\pi_i(1)}_{\varepsilon_i}\Delta^{\pi_i(2)}_{\varepsilon_i}\cdots\Delta^{\pi_i(l_i)}_{\varepsilon_i}}_{\Delta_{\varepsilon_i}^{\Sigma_j\pi_i(j)\overset{\text{def}}{=}\mathrm{long}(\pi_i)}}, % [sic] \Delta_{\varepsilon_1} -> \Delta_{\varepsilon_i}, and g_{\varepsilon,\pi_i} -> g_{\varepsilon_i,\pi_i}, index i intended
\tag*{(1)}
\]

\msp{page 53}
where $g_{\varepsilon_i,\pi_i}$ is the arrow defined in terms of $\pi_i$ in the preceding paragraph. But this is also nothing other than
\[
g_{\varepsilon_i,\pi_i} = \mathfrak{f}_{\varepsilon_i,\pi_i(1)}*\mathfrak{f}_{\varepsilon_i,\pi_i(2)}*\cdots*\mathfrak{f}_{\varepsilon_i,\pi_i(l_i)}.
\tag*{(2)}
\]
The element of $\mathrm{Nerf}_2$ under consideration depends only on
\[
(\varepsilon_i,l_i,\pi_i), \qquad\quad \text{where } (\varepsilon_i,l_i)\in\{\pm1\}\times\mathbf{N},\ \pi_i : [1,l_i]\longrightarrow\mathbf{N}.
\]
(\textbf{NB} If $l_i=0$, $[1,l_i]=[1,0]=\emptyset$ by convention, and $\pi_i$ is the unique map $\emptyset\longrightarrow\mathbf{N}$, so the datum $(\varepsilon_i,l_i,\pi_i)$ then amounts to $(\varepsilon_i,l_i)$.)

If
\[
(\varepsilon,l,\pi)\ \text{ as above,} \qquad\quad (\varepsilon,l)\in\{\pm1\}\times\mathbf{N},\ \pi : [1,l]\longrightarrow\mathbf{N},
\]
let us denote by
\[
\mathfrak{f}_{\varepsilon,l,\pi}
\]
the element of $\mathrm{Nerf}_2$ given by
\[
\Delta_\varepsilon\xrightarrow{\mathfrak{f}_{\varepsilon,l}}\Delta^l_\varepsilon\ \xrightarrow{\quad\mathfrak{f}_{\varepsilon,\pi(1)}\ \ *\ \cdots\ *\ \ \mathfrak{f}_{\varepsilon,\pi(l)}\quad}\Delta^{\pi(1)}_\varepsilon\Delta^{\pi(2)}_\varepsilon\cdots\Delta^{\pi(l)}_\varepsilon.
\]
We can formulate the preceding construction by saying that every element of $\mathrm{Nerf}_2$ can be decomposed uniquely into a product
\[
(f,g) = \mathfrak{f}_{\varepsilon_1,l_1,\pi_1}*\mathfrak{f}_{\varepsilon_2,l_2,\pi_2}*\cdots*\mathfrak{f}_{\varepsilon_n,l_n\pi_n}, % [sic] l_n\pi_n -> l_n,\pi_n, comma missing in index
\]
which proves that $\mathrm{Nerf}_2(\uTy)$ is the free monoid on the preceding generators $\mathfrak{f}_{\varepsilon,l,\pi}$. (\textbf{NB} to the empty sequence corresponds
\[
1\xrightarrow{\mathrm{id}_1}1\xrightarrow{\mathrm{id}_1}1 \qquad\qquad .)
\]
% foot: Version 22 décembre 2001

\printedpage{37}
% head: A. Grothendieck : Dérivateurs, Ch. VII, Édition des Universités de Montpellier II et Paris VII
\msp{page 54}
If the object of $\mathrm{Nerf}_2$ is
\[
x\xrightarrow{\ f\ }y\xrightarrow{\ g\ }z,
\]
the preceding construction tells us, in terms of the combinatorial datum, i.e. of the sequence of data
\[
(\varepsilon_1,l_1,\pi_1),\ \dots\ ,\ (\varepsilon_n,l_n,\pi_n),
\]
what $x$, $y$, $z$ are: each of them is the product in $\mathrm{Nerf}_0=\mathrm{Ty}$ of the $n$ factors $x_i$, $y_i$ or $z_i$ respectively, for $i\in[1,n]$, given in the definition of the $\mathfrak{f}_{\varepsilon_i,l_i,\pi_i}$, i.e. of
\[
\Delta_{\varepsilon_i},\ \Delta^{l_i}_{\varepsilon_i},\ \text{ and }\Delta^{\pi_i(1)}_{\varepsilon_i}\Delta^{\pi_i(2)}_{\varepsilon_i}\cdots\Delta^{\pi_i(l_i)}_{\varepsilon_i}=\Delta^{\mathrm{long}(\pi_i)}_{\varepsilon_i}
\]
(where $\mathrm{long}(\pi_i)=\sum_{j\in[1,l_i]}\pi_i(j)$). Thus
\[
\left\{
\begin{array}{lcl}
x & = & \prod^{\mathrm{Ty}}_{i\in[1,n]}\Delta_{\varepsilon_i}\\[4pt]
y & = & \prod^{\mathrm{Ty}}_{i\in[1,n]}\Delta^{l_i}_{\varepsilon_i}\\[4pt]
z & = & \prod^{\mathrm{Ty}}_{i\in[1,n]}\Delta^{\mathrm{long}(\pi_i)}_{\varepsilon_i}
\end{array}
\right.
\]
(\footnote[16]{careful, the products must be taken respecting the order of the factors given by the index notation in $i$.}), and we similarly find $f$, $g$ as products (in the sense $*$) of $n$ factors $f_i$ resp. $g_i$ ($i\in[1,n]$) made explicit in formulas (1), (2) (p. 52, 53):
\[
\begin{array}{lclccl}
f_i & = & \mathfrak{f}_{\varepsilon_i,l_i} : & \Delta_{\varepsilon_i} & \longrightarrow & \Delta^{l_i}_{\varepsilon_i}\\[4pt]
g_i & = & g_{\varepsilon_i,l_i} : & \Delta^{l_i}_{\varepsilon_i} & \longrightarrow & \Delta^{\mathrm{long}(\pi_i)}_{\varepsilon_i}, % [sic] g_{\varepsilon_i,l_i} -> g_{\varepsilon_i,\pi_i}, notation of (1), (2)
\end{array}
\]
$g_i$ itself being represented as a $*$-product of $l_i$ indecomposable factors, so that $f$ and $g$ appear in the form
\msp{page 55}
of $*$-products of $n$, resp. of $\sum_{i\in[1,n]}l_i$ factors of the form $\mathfrak{f}_{\varepsilon,\nu}$. To finish getting the truncated semi-simplicial object $(\mathrm{Nerf}_i(\uTy))_{0\leq i\leq 2}$ in hand, it remains to describe the two degeneracy maps $\mathrm{Nerf}_1\longrightarrow\mathrm{Nerf}_2$ (the case where $f$ or $g$ is an identity arrow) and above all the \textit{composition} map for arrows
\[
\mathrm{Nerf}_2(\uTy)\xrightarrow{\ \gamma\ }\mathrm{Nerf}_1(\uTy)=\mathrm{Fl}(\uTy).
\]
It suffices to give $\gamma$ on the generators, and we have tautologically, by definition of the $\mathfrak{f}_{\varepsilon,l,\pi}$
\[
\gamma(\mathfrak{f}_{\varepsilon,l,\pi})=\mathfrak{f}_{\varepsilon,\mathrm{long}(\pi)}.
\]
(\textbf{NB} There is at most one morphism from $\Delta_\varepsilon$ (from $\Delta_+$ or from $\Delta_-$) to an interval $I_\tau$. It exists if and only if $\tau$ is of the form $\Delta^l_\varepsilon$, $l\in\mathbf{N}$.)
% foot: Version 22 décembre 2001

\printedpage{38}
% head: A. Grothendieck : Dérivateurs, Ch. VII, Édition des Universités de Montpellier II et Paris VII
The case $g=\mathrm{id}$ corresponds to the case where
\[
\pi_i(j)=1 \qquad\quad \forall j\in[1,l_i],
\]
the $\mathfrak{f}_{\varepsilon_i,l_i,\pi_i}$ are then determined in terms of the $\mathfrak{f}_{\varepsilon_i,l_i}$, and the degeneracy map is % [sic] détérminés -> déterminés, accent
\[
\mathfrak{f}_{\varepsilon,l}\longmapsto\mathfrak{f}_{\varepsilon,l,1_l} \qquad\quad (\text{where } 1_l \text{ is the constant function with value 1 on } [1,l]).
\]
The $f=\mathrm{id}$ is the one where the $l_i$ are equal to 1. Then each $\pi_i : [1,1]\longrightarrow\mathbf{N}$ can be viewed as an element of $\mathbf{N}$, giving $(\varepsilon_i,l_i,\pi_i)$ amounts to giving the $(\varepsilon_i,n_i)\in\{\pm1\}\times\mathbf{N}$, i.e. giving % [sic] Le f = id -> Le cas f = id, word « cas » missing
\msp{page 56}
a sequence of $n$ generators among the $\mathfrak{f}_{\varepsilon,l}$. The degeneracy map $\mathrm{Nerf}_1\longrightarrow\mathrm{Nerf}_2$ is here given on the generators by
\[
\mathfrak{f}_{\varepsilon,l}\longmapsto\mathfrak{f}_{\varepsilon,1,l}.
\]
Fortunately, we no longer have to prove the associativity of composition.

\textbf{NB} The simplicity of these descriptions reinforces my feeling that the category $\uTy$ now defined, with the \textit{increasing morphisms of oriented intervals} $I_\tau\longrightarrow I_{\tau'}$ as arrows of $\uTy$, is indeed ``the right one''. We shall see for a start whether it gives a nice description of the categories of fractions $X\Sigma^{-1}$, without contortions as before.

But we have promised to determine $\mathrm{Nerf}_r(\uTy)$, for every $r\in\mathbf{N}$. The general generator of $\mathrm{Nerf}_r$ is of the form (for $\varepsilon\in\{\pm1\}$ fixed)
\[
\underbrace{\Delta_\varepsilon}_{=\,\Delta^{l_0}_\varepsilon\text{with } l_0=1}\xrightarrow{\ f_1\ }\Delta^{l_1}_\varepsilon\xrightarrow{\ f_2\ }\Delta^{l_2}_\varepsilon\longrightarrow\ \cdots\ \xrightarrow{\ f_r\ }\Delta^{l_r}_\varepsilon,
\]
where the $f_i : \Delta^{l_{i-1}}_\varepsilon\longrightarrow\Delta^{l_i}_\varepsilon$, $i\in[1,r]$, are given by maps
\[
\pi_{i-1} : [1,l_{i-1}]\longrightarrow\mathbf{N} \qquad\quad \text{such that }\ \underbrace{\mathrm{long}(\pi_{i-1})}_{\overset{\text{def}}{=}\ \Sigma_{j\in[1,l_i]}\pi_i(j)}\ = l_i. % [sic] \Sigma_{j\in[1,l_i]}\pi_i(j) -> \Sigma_{j\in[1,l_{i-1}]}\pi_{i-1}(j), indices shifted with respect to long(\pi_{i-1})
\]
(NB $\pi_0 : [1,1]=\{1\}\longrightarrow\mathbf{N}$ is identified with $\mathrm{long}(\pi_0)=\pi_0(1)=l_1$.)

\msp{page 57}
To avoid making the definition of the domain of the $\pi_i$ ($i\in[0,r-1]$) depend on the integers $l_i$, we can also consider the $\pi_i$ as maps
\[
\mathbf{N}^*\xrightarrow{\ \pi\ }\mathbf{N} \qquad\quad (\mathbf{N}^*=\{i\in\mathbf{N}\mid i\geq 1\})
\]
with finite supports (i.e. such that $\pi(i)=0$ for $i$ large), and require precisely
\[
\begin{aligned}
&\mathrm{supp}\,\pi_0\ \subseteq\ \{1\}\\
&\mathrm{supp}\,\pi_1\ \subseteq\ [1,\mathrm{long}(\pi_0)]\\
&\qquad\ \cdots\\
&\boxed{\mathrm{supp}\,\pi_i\ \subseteq\ [1,\mathrm{long}(\pi_{i-1})]}, \qquad\quad \text{for }\ \boxed{i\in[0,r-1]}\quad .
\end{aligned}
\tag*{$(**)$}
\]
% foot: Version 22 décembre 2001

\end{document}
